\documentclass{article}
\usepackage{amsmath, amssymb, amsthm}

\theoremstyle{plain}
\newtheorem{theorem}{Theorem}
\newtheorem{lemma}[theorem]{Lemma}
\newtheorem*{answer}{Answer}

\begin{document}

\title{Non-Existence of Compact Manifolds with Rationally Acyclic Universal Covers for Lattices with Torsion}
\author{}
\date{}
\maketitle

\begin{abstract}
We address the question of whether a uniform lattice  in a real semi-simple Lie group, which contains 2-torsion, can be the fundamental group of a compact manifold without boundary whose universal cover is acyclic over the rational numbers . Using the Hattori-Stallings trace on the Grothendieck group , we prove that the existence of such a manifold is impossible due to the presence of torsion.
\end{abstract}

\section{Answer}

\begin{answer}
No, it is not possible.
\end{answer}

\section{Mathematical Reasoning}

We establish the non-existence of such a manifold by deriving a contradiction from the existence of a non-trivial torsion element in  and the finite dimensionality of the manifold .

Let  be a uniform lattice in a real semi-simple Lie group, and suppose  contains an element  of order 2. Assume, for the sake of contradiction, that there exists a compact manifold  without boundary such that  and the universal covering space  is acyclic over .

\subsection{Algebraic Formulation}

Since  is a compact manifold, it has the homotopy type of a finite CW-complex. Let . The cellular chain complex of the universal cover with rational coefficients, denoted , consists of free -modules. Since  is finite, these modules are finitely generated. Specifically, if  has  cells in dimension , then

The condition that  is -acyclic means that its reduced rational homology vanishes. Consequently, the augmented chain complex is exact:
\begin{equation}
0 \to C_n(\tilde{M}; \mathbb{Q}) \to \dots \to C_0(\tilde{M}; \mathbb{Q}) \to \mathbb{Q} \to 0,
\end{equation}
where  denotes the trivial -module.

This exact sequence implies an equality in the Grothendieck group  of finitely generated projective -modules:
\begin{equation}
[\mathbb{Q}] = \sum_{k=0}^n (-1)^k [C_k(\tilde{M}; \mathbb{Q})].
\end{equation}

\subsection{The Hattori-Stallings Trace}

We employ the Hattori-Stallings trace (or character trace), a homomorphism , where  is the set of conjugacy classes of . For a group element , let  denote the coefficient of the conjugacy class of  in the trace of a projective module . Equivalently, this is the character of the representation  evaluated at .

We evaluate the trace at the torsion element .

\begin{enumerate}
\item \textbf{Left-Hand Side:} The module  is the trivial representation. The action of any element  is the identity map on . Therefore, the character value is:

```
\item \textbf{Right-Hand Side:} Consider the term $[C_k(\tilde{M}; \mathbb{Q})]$. This is the class of a free module $(\mathbb{Q}\Gamma)^{c_k}$. The action of $\gamma$ on the basis $\Gamma$ is given by right multiplication $x \mapsto x\gamma$. Since $\Gamma$ acts freely on itself, and $\gamma \neq 1$, this action permutes the basis elements without fixed points. In the regular representation, the trace of any non-identity element is zero. Thus:
\[
T_\gamma([C_k(\tilde{M}; \mathbb{Q})]) = 0.
\]

```

\end{enumerate}

\subsection{Contradiction}

Applying the trace  to the equality in equation (2), we obtain:

The equation  is a contradiction.

\section{Conclusion}

The contradiction arises from the assumption that a group with non-trivial torsion can act freely on a finite-dimensional -acyclic complex (specifically, the universal cover of a compact manifold). Therefore,  cannot be the fundamental group of a compact manifold whose universal cover is acyclic over .

\end{document}
